% --- IMPORTS ---
\documentclass[leqno]{article}
\usepackage{graphicx}
\usepackage{dsfont}
\usepackage{mathtools}
\usepackage{amssymb}
\usepackage{amsmath}
\usepackage{amsthm}
\usepackage{float}
\usepackage{ragged2e}
\usepackage{tensor}
\usepackage{fancyhdr}
\usepackage{empheq}
\usepackage[most]{tcolorbox}
\usepackage{enumitem}
\usepackage{dirtytalk}
\usepackage{marginnote}
\usepackage{microtype}
\usepackage[
  a4paper,
  left=0.8in,
  right=1in,
  top=1in,
  bottom=0.5in,
  headheight=14pt,
  headsep=0.4in
]{geometry}
\setlength{\parindent}{0cm}
% --- TikZ Packages ---
\usepackage{pgfplots}
\pgfplotsset{compat=1.18}
\usepgfplotslibrary{fillbetween, polar}
\usetikzlibrary{calc, positioning}
\usetikzlibrary{angles, quotes}
\usetikzlibrary{arrows.meta}
\usetikzlibrary{decorations.pathreplacing, decorations.markings}
\usetikzlibrary{patterns}
\usetikzlibrary{intersections}
% --- URL Packages ---
\usepackage[colorlinks=true,linkcolor=red,citecolor=magenta,urlcolor=black]{hyperref}%
\urlstyle{same}
% --- END OF IMPORTS ---

% --- CUSTOM COMMANDS ---
\RenewDocumentCommand{\marks}{o m}{%
  \IfValueTF{#1}%
    {\marginnote{\raggedright\textbf{[#2]}}[#1]}%
    {\marginnote{\raggedright\textbf{[#2]}}}%
}
\newcommand{\vspacerule}[2]{%
  \par\vspace*{#1}%
  \noindent\hrulefill\par
  \vspace*{#2}%
}
\definecolor{scarlet}{RGB}{205, 40, 75}
\newcommand{\questionitem}{%
  \item\phantomsection
  \addcontentsline{toc}{section}{Question \number\value{enumi}}%
}
\renewcommand{\contentsname}{Questions}
% --- END OF CUSTOM COMMANDS ---

% --- HEADER CONFIG ---
\pagestyle{fancy}
\fancyhead{}
\fancyfoot{}
\renewcommand{\headrulewidth}{0.9pt}
\lhead{Complex Series}
\chead{}
\rhead{\href{https://danieljsmith.org}{danieljsmith.org}}
% --- END OF HEADER CONFIG ---

\begin{document}
\vspace*{20pt}
\tableofcontents
\newpage
\begin{enumerate}[
  label=\textbf{\arabic*.},
  leftmargin=*,
  topsep=0pt,
  partopsep=0pt,
  parsep=0pt
]

% ---- Question 1 ----
\questionitem

The infinite series $C$ and $S$ are defined by
\begin{align*}
 C&= \cos 2\theta - \frac{1}{3}\cos 5\theta + \frac{1}{9}\cos 8\theta - \frac{1}{27}\cos 11\theta + \cdots \\[4pt]
S &= \sin 2\theta - \frac{1}{3}\sin 5\theta + \frac{1}{9}\sin 8\theta - \frac{1}{27}\sin 11\theta + \cdots  
\end{align*}

Given that the series $C$ and $S$ are both convergent. \\

\begin{enumerate}[label=\textbf{(\alph*)}]
  \item Show that
  \[
  C+ \mathrm{i}S = \frac{3e^{2\mathrm{i}\theta}}{3 + e^{3\mathrm{i}\theta}}
  \]
  \marks[-20pt]{4}
  \item Hence show that
  \[
  S = \frac{9\sin 2\theta - 3\sin \theta}{10 + 6\cos 3\theta}
  \]
  \marks[-30pt]{4}
\end{enumerate}

\newpage


% ---- Question 2 ----
\questionitem

\begin{enumerate}[label=(\roman*)]
  \item Show that \[(1 - 2e^{2\mathrm{i}\theta})(1 - 2e^{-2\mathrm{i}\theta}) = 5 - 4\cos 2\theta\] \marks[-15pt]{3} \\
\end{enumerate}

For a positive integer $n$, series $C$ and $S$ are defined by\\

\[
\begin{aligned}
C &= \cos\theta + 2\cos 3\theta + 4\cos 5\theta + \cdots + 2^{n-1}\cos((2n-1)\theta)\\[5pt]
S &= \sin\theta + 2\sin 3\theta + 4\sin 5\theta + \cdots + 2^{n-1}\sin((2n-1)\theta)
\end{aligned}
\] \\
\begin{enumerate}[label=(\roman*), resume]
  \item Show that
  \[
  C = \frac{2^{n+1}\cos((2n-1)\theta) - 2^n\cos((2n+1)\theta) - \cos\theta}{5 - 4\cos 2\theta}
  \]
  and find a similar expression for $S$. \marks{9}
\end{enumerate}

\newpage

% ---- Question 3 ----
\questionitem

\begin{enumerate}[label=\textbf{(\roman*)}]
  \item Show that
  \[
  1 - e^{i2\theta} = 2\sin\theta(\sin\theta - i\cos\theta)
  \] \marks[-15pt]{3}\\
  \item For a positive integer $n$, the series $C$ and $S$ are defined as follows.
  \begin{align*}
  C &= 1 - \binom{n}{1}\cos 2\theta + \binom{n}{2}\cos 4\theta - \ldots + (-1)^n\cos 2n\theta \\[5pt]
  S &=\quad -\binom{n}{1}\sin 2\theta + \binom{n}{2}\sin 4\theta - \ldots + (-1)^n\sin 2n\theta
  \end{align*} \\
  By considering $C + iS$, show that
  \[
  C = 2^n\sin^n\theta\cos\left(n\theta - \frac{n\pi}{2}\right)
  \]
  and find a corresponding expression for $S$.\marks{9} \\
\end{enumerate}

\newpage

% ---- Question 4 ----
\questionitem

The infinite series $C$ and $S$ are defined as follows: \\
\begin{align*}
C&= \cos\alpha + r\cos(\alpha + \beta) + r^2\cos(\alpha + 2\beta) + \ldots \\[10pt]
S &= \sin\alpha + r\sin(\alpha + \beta) + r^2\sin(\alpha + 2\beta) + \ldots \\
\end{align*}

where $r$ is a real number and $|r| < 1$. \\


By considering $C+ \mathrm{i}S$, show that \\
\[
C= \frac{\cos\alpha - r\cos(\alpha - \beta)}{1 + r^2 - 2r\cos\beta} \\
\] 
and find a corresponding expression for $S$. \marks{8}

\newpage

% ---- Question 5 ----
\questionitem

\begin{enumerate}[label=(\roman*)]
  \item For values of $\theta$ for which the expressions are defined, show that
  \[
  1+\mathrm{i}\tan\theta=\sec\theta(\cos\theta+\mathrm{i}\sin\theta)
  \]
  \marks[-20pt]{2}

  \item For a positive integer $n$, the series $C$ and $S$ are defined by
  \begin{align*}
    C&=\,1-\binom{n}{2}\tan^2\theta+\binom{n}{4}\tan^4\theta-\cdots\\[10pt]
  S&=\qquad\binom{n}{1}\tan\theta-\binom{n}{3}\tan^3\theta+\binom{n}{5}\tan^5\theta-\cdots
  \end{align*}

  where in each case the pattern continues until the power of $\tan\theta$ is at most $n$. \\

  By considering $C+\mathrm{i}S$, show that
  \[
  C=\sec^n\theta\cos n\theta
  \]
  and find a corresponding expression for $S$. \marks{7}
\end{enumerate}
\newpage

% ---- Question 6 ----
\questionitem

The convergent series $C$ and $S$ are defined by
\begin{align*}
C &= \sum_{n=0}^{\infty} \left(-\frac{1}{2}\right)^n \cos\bigl((4n+1)\theta\bigr)
\\[10pt]
S &= \sum_{n=0}^{\infty} \left(-\frac{1}{2}\right)^n \sin\bigl((4n+1)\theta\bigr)
\end{align*} \\
\begin{enumerate}[label=\textbf{(\alph*)}]
  \item Show that
  \[
  C + \mathrm{i}S = \frac{2e^{\mathrm{i}\theta}}{2 + e^{4\mathrm{i}\theta}}
  \]
  \marks[-20pt]{4}
  \item Hence show that
  \[
  C = \frac{4\cos \theta + 2\cos 3\theta}{5 + 4\cos 4\theta}
  \]
  \marks[-30pt]{4}
\end{enumerate}
\newpage
% ---- Question 7 ----
\questionitem

Let
\begin{align*}
C&= \sum_{r=0}^{n-1}\cos^2\left(\theta + \frac{2\pi r}{n}\right)\\[10pt]
S &= \sum_{r=0}^{n-1}\sin\left(\theta + \frac{2\pi r}{n}\right)\cos\left(\theta + \frac{2\pi r}{n}\right)\\
\end{align*}

where $n$ is an integer greater than $2$. Let $\exp(x)$ denote $e^x$.\\


By considering the real and imaginary parts of
\[
\sum_{r=0}^{n-1} \exp\left(2\mathrm{i}\left(\theta + \dfrac{2\pi r}{n}\right)\right)
\]
show that $C= \dfrac{n}{2}$ and $S = 0$. \marks{7} \\

\newpage
\end{enumerate}
\end{document}
